\documentclass[tikz,border=4pt]{standalone}
% Paleta y fuentes compartidas por el libro y las figuras (única fuente de verdad).
\RequirePackage{xcolor}
\definecolor{oxido}{HTML}{B7410E}
\definecolor{cobre}{HTML}{D27D2D}
\definecolor{carbon}{HTML}{2B2B2B}
\definecolor{azul}{HTML}{1F5F8B}
\definecolor{arena}{HTML}{F4EDE4}

\usepackage{fontspec}
\setsansfont{FiraSans}[Extension=.otf,UprightFont=*-Regular,BoldFont=*-Bold,ItalicFont=*-Italic,BoldItalicFont=*-BoldItalic]
\setmonofont{FiraMono}[Extension=.otf,UprightFont=*-Regular,BoldFont=*-Medium]
\usetikzlibrary{positioning,arrows.meta,calc}
\begin{document}
\sffamily\small
\begin{tikzpicture}[
  celda/.style={draw=carbon,minimum height=7mm,minimum width=14mm,font=\ttfamily\small},
  byte/.style={draw=carbon,minimum height=7mm,minimum width=7mm,font=\ttfamily\small,fill=white},
  rotulo/.style={font=\sffamily\footnotesize\color{carbon!75}},
  flecha/.style={->,very thick,rounded corners=6pt},
  >={Stealth[length=2.6mm]}]
  % Pila: el String y dos rebanadas
  \node[font=\sffamily\bfseries] at (1.4,1.2) {Pila (\textit{stack})};
  \foreach \c/\v [count=\i from 0] in {ptr/{}, len/10, cap/10} {
    \node[celda,fill=arena] (s\i) at (\i*1.4,0) {\v};
    \node[rotulo] at (\i*1.4,0.55) {\c};
  }
  \node[anchor=east,font=\ttfamily\bfseries] at (-0.9,0) {s};
  \foreach \c [count=\i from 0] in {ptr,len} {
    \node[celda,fill=azul!12] (p\i) at (\i*1.4,-2.0) {};
    \node[rotulo] at (\i*1.4,-1.45) {\c};
    \node[celda,fill=azul!12] (m\i) at (\i*1.4,-3.8) {};
    \node[rotulo] at (\i*1.4,-3.25) {\c};
  }
  \node[font=\ttfamily\small] at (p1) {4};
  \node[font=\ttfamily\small] at (m1) {5};
  \node[anchor=east,font=\ttfamily\bfseries] at (-0.9,-2.0) {hola};
  \node[anchor=east,font=\ttfamily\bfseries] at (-0.9,-3.8) {mundo};
  % Montículo: los 10 bytes de "hola mundo"
  \node[font=\sffamily\bfseries] at (8.65,-0.75) {Montículo (\textit{heap})};
  \foreach \l [count=\i from 0] in {h,o,l,a,{\color{carbon!45}\rule[-1pt]{3mm}{0.6pt}},m,u,n,d,o} {
    \node[byte] (h\i) at (5.5+\i*0.7,-2.0) {\l};
    \node[rotulo] at (5.5+\i*0.7,-1.45) {\i};
  }
  % Punteros: salen hacia abajo y van por pasillos libres
  \fill[carbon] (s0.center) circle (1.6pt);
  \draw[flecha,carbon] (s0.center) -- (0,-0.95) -- (4.55,-0.95)
    -- (4.55,-1.85) -- ([yshift=1.5mm]h0.west);
  \fill[azul] (p0.center) circle (1.6pt);
  \draw[flecha,azul] (p0.center) -- (0,-2.8) -- (4.55,-2.8)
    -- (4.55,-2.15) -- ([yshift=-1.5mm]h0.west);
  \fill[azul] (m0.center) circle (1.6pt);
  \draw[flecha,azul] (m0.center) -- (0,-4.6) -- (h5.south |- 0,-4.6)
    -- (h5.south);
  % Código y nota
  \node[anchor=west,font=\ttfamily\small,align=left,fill=fondocodigo,inner sep=5pt]
    at (-1.6,-5.75) {let s = String::from("hola mundo");\\let hola = \&s[..4];\\let mundo = \&s[5..];};
  \node[anchor=west,font=\sffamily\footnotesize,text=carbon!75,align=left]
    at (6.0,-5.75) {Una rebanada son dos números:\\dónde empieza y cuántos bytes mide.\\No copia el texto ni es su dueña.};
\end{tikzpicture}
\end{document}
